% TOP_PROBLEM: {"id": 2800604, "title": "10 Lectures and 42 Open Problems — The Paley ETF Conjecture", "category_id": 15, "category": {"id": 15, "name": "computer_science", "display_name": "Computer Science", "description": "Computational complexity, algorithms, and theoretical CS.", "slug": "computer-science", "order_index": 15, "created_at": "2026-07-31T15:26:25.671Z"}, "status": "partially_solved", "problem_number": "AMR-027-0604", "set_id": 13}
% FIRST_POSTED: 2026-10-02
% ENTRY_KIND: statement_only
% UnsolvedMath ID 2800604; source catalogue code AMR-027-0604
% !TeX program = lualatex
% Definition and sources only; this file is not an LLM solution attempt.
\documentclass[11pt,a4paper]{article}
\usepackage[margin=25mm]{geometry}
\usepackage{fontspec}
\setmainfont{lmroman10-regular.otf}[BoldFont=lmroman10-bold.otf,
 ItalicFont=lmroman10-italic.otf,BoldItalicFont=lmroman10-bolditalic.otf]
\usepackage{amsmath,amssymb,mathtools}
\usepackage{unicode-math}
\setmathfont{latinmodern-math.otf}
\AtBeginDocument{\let\setminus\smallsetminus}
\usepackage{xurl,hyperref}
\hypersetup{colorlinks=true,urlcolor=blue}
\setcounter{secnumdepth}{0}
\setlength{\parindent}{0pt}
\setlength{\parskip}{5pt}
\setlength{\emergencystretch}{3em}
\begin{document}
\section*{10 Lectures and 42 Open Problems — The Paley ETF Conjecture}\label{2800604}
{\small UnsolvedMath ID 2800604 · AMR-027-0604 · Computer Science · AMR Open Problem Lists}
\subsection{Definitions and mathematical statement}
Let $p\equiv1\pmod4$ be prime and let $Q\subset\mathbb F_p^*$ be
the nonzero quadratic residues. Form a matrix $A_p$ with rows
$\{0\}\cup Q$ and columns $\mathbb F_p\cup\{\infty\}$ by
\[
 (A_p)_{0j}=p^{-1/2},\quad
 (A_p)_{rj}=(2/p)^{1/2}e^{2\pi i rj/p}\quad(r\in Q,j\in\mathbb F_p),
 \qquad (A_p)_{\cdot,\infty}=e_0.
\]
The vector $e_0$ is one in row zero and zero elsewhere. This is the
unit-column-norm Paley equiangular tight frame, of size
$(p+1)/2$ by $p+1$. Distinct columns have inner-product magnitude
$p^{-1/2}$.

For $s\ge1$, set
\[
 \delta_s(A_p)=\max_{|J|\le s}
                 \|(A_p)_J^*(A_p)_J-I_{|J|}\|_{\mathrm{op}},
\]
where $J$ is a column subset and the norm is the largest singular
value. Does there exist a fixed $\delta<1$ and an integer function
$s(p)$ with $s(p)/\sqrt p\to\infty$ such that
$\delta_{s(p)}(A_p)\le\delta$ for all sufficiently large such primes?

Even an unbounded logarithmic improvement over $\sqrt p$ is sought;
a fixed constant multiple of $\sqrt p$ does not suffice. Equivalently,
all coefficient vectors on each permissible column set must have their
squared norm preserved within factors $1\pm\delta$. The question is
unconditional for this specific arithmetic frame, not a random choice
of Fourier rows or a conclusion assuming an unproved distribution
property of quadratic residues.
\subsection{Short English statement}
Does the Paley frame preserve sparse-vector norms for sparsity growing faster than the square root of its size?
\subsection{Sources}
\begin{itemize}
\item[S1] ULAM.AI and the UnsolvedMath Contributors, supplied problems.json, record 2800604 (AMR-027-0604), AMR Open Problem Lists. Catalogue identity and source transcription; CC BY 4.0.
\url{https://huggingface.co/datasets/ulamai/UnsolvedMath}.
\item[S2] Bandeira - 42 Open Problems in Data Science (2016), Open Problem 6.4. Original source indexed by AMR-027-0604.
\url{https://afonsobandeira.wordpress.com/2015/12/07/10l42pcompressedsensing/}.
\item[S3] A. S. Bandeira, Ten Lectures and Forty-Two Open Problems in the Mathematics of Data Science, corresponding numbered problem and surrounding definitions.
\url{https://people.math.ethz.ch/~abandeira/TenLecturesFortyTwoProblems.pdf}.
\end{itemize}
\end{document}
